% apore-lite condensed study notes.
% Generated from this chapter's wiki/. Safe to edit by hand.
% Keep both APORE BODY markers exactly as they are:
% shared/scripts/build_notes.py extracts the text between them when
% building the combined domain PDF.
\documentclass[11pt]{article}
\usepackage{apore-notes}

\notesmeta{Fundamentals of Engineering Materials}{1. Introduction}{2026-09-19}

\begin{document}
\makenotestitle
% >>> APORE BODY START
\chapterhead{Introduction}

\topic{Materials Science and Engineering}

\begin{keydef}{Materials Science}
Investigates relationships between structures and properties of materials, and
designs/develops new materials. \src{2.PNG}
\end{keydef}

\begin{keydef}{Materials Engineering}
Creates products from existing materials and develops materials processing
techniques. \src{2.PNG}
\end{keydef}

\begin{keydef}{Core terms}
\begin{points}
\item \textbf{Composition}: chemical make-up of a material.
\item \textbf{Structure}: arrangement of atoms in a material.
\item \textbf{Processing}: ways to shape materials into useful components.
\end{points}
\src{2.PNG}
\end{keydef}

\begin{formula}{The materials chain}
\[
\text{Composition} \to \text{Structure} \to \text{Processing} \to
\text{Properties} \to \text{Performance}
\]
\src{2.PNG}
\end{formula}

\begin{keydef}{Materials knowledge spectrum}
Materials science provides \emph{basic} knowledge; materials engineering
provides \emph{applied} knowledge. Together they produce the resultant knowledge
of structure, properties, processing and performance that engineers use to
convert materials into the products needed by society. \src{6.PNG}
\end{keydef}

\begin{worked}{Structure, not composition, sets properties}
A disk of aluminum oxide in three forms (single crystal, polycrystalline, and
polycrystalline with many pores) has different optical properties in each form
despite identical composition. \src{3.PNG}
\end{worked}

\topic{Why Materials Are Important}

\begin{keydef}{Why it matters}
Materials drive advancements in society. Engineers must understand materials in
order to select appropriate materials and processing techniques for specific
applications. \src{4.PNG, 5.PNG}
\end{keydef}

\begin{keydef}{Historical materials ages}
\begin{points}
\item Stone Age, Bronze Age, Iron Age mark society's advancement.
\item Today's age is posed as an open question. The candidates offered are the
Silicon (Electronic Materials) Age, the Nano-materials Age, and the Polymer Age.
\end{points}
\src{4.PNG}
\end{keydef}

\begin{keydef}{Engineer's requirement}
Products, devices and components are all made of materials. Selecting
appropriate materials and processing requires knowledge of material properties
and of structure--property relationships. \src{5.PNG}
\end{keydef}

\topic{Types of Materials}

\begin{keydef}{Three broad classes}
Metals and alloys; ceramics and glasses; polymers and plastics. \src{7.PNG}
\end{keydef}

\subtopic{Metals and Alloys}

\begin{keydef}{Alloy}
A metal that contains additions of one or more metals or non-metals.
\src{8.PNG}
\end{keydef}

\begin{keydef}{Metals: structure and properties}
\begin{points}
\item Pure metallic elements or combinations of them; a large number of
de-localized electrons, and therefore able to conduct electricity. \src{7.PNG}
\item Properties: good thermal and electrical conductivity, high strength, high
stiffness, ductility. \src{8.PNG}
\item Steel is an alloy of iron and a small content of carbon. Brass is an alloy
of copper and zinc. \src{8.PNG}
\item Examples: aluminum, magnesium, zinc, iron, titanium, beryllium-copper,
nickel, copper-nickel, niobium. \src{8.PNG}
\end{points}
\end{keydef}

\subtopic{Ceramics and Glasses}

\begin{keydef}{Ceramics}
Compounds of metallic and non-metallic elements: oxides, carbides, nitrides,
sulfides; generally insulating and refractory. \src{7.PNG} Inorganic
non-metallic materials, crystalline (polycrystalline) in structure; examples
include sand and rocks. Properties: insulating (non-conductive), high hardness,
high strength and melting point, very brittle. \src{9.PNG}
\end{keydef}

\begin{keydef}{Glasses}
Amorphous materials with structure-dependent properties, which can be processed
to yield an improvement in properties. \src{9.PNG}
\end{keydef}

\subtopic{Polymers and Plastics}

\begin{keydef}{Polymer}
``Poly'' means many, ``mer'' means unit. Polymers are large molecular structures
formed from monomers through polymerization. \src{10.PNG}
\end{keydef}

\begin{keydef}{Polymers: composition and properties}
\begin{points}
\item Chemical composition including carbon, hydrogen and other non-metals;
large molecular structure; flexible, lightweight, easily processed. \src{7.PNG}
\item Synthetic: polyethylene (PE), polymethyl methacrylate (PMMA),
polycarbonate. Natural: cellulose, glucan. \src{10.PNG}
\end{points}
\end{keydef}

\begin{points}
\item \textbf{Practice.} Classify as metal, ceramic or polymer: brass,
$\mathrm{MgO}$, plexiglass, polychloroprene, $\mathrm{B_4C}$, cast iron. The
source provides no answer key. \src{14.PNG}
\end{points}

\topic{Other Materials}

\begin{keydef}{Composite}
A combination of two or more different materials, combining the properties of
the materials it is formed from. \src{11.PNG, 12.PNG}
\end{keydef}

\begin{keydef}{Advanced materials}
Semiconductors, biomaterials, smart materials, and nanoengineered materials.
\src{11.PNG}
\end{keydef}

\begin{keydef}{Semiconductors}
Conductivity in between ceramics and metals. Used in computers, electronics and
micro-electrical-mechanical systems (MEMS). Examples: silicon, germanium,
gallium arsenide. \src{13.PNG}
\end{keydef}

\topic{Comparing Material Properties}

\begin{keydef}{Material property chart}
Compares classes of materials (metals, ceramics, polymers, composites) against
measurable properties such as density, elastic modulus and tensile strength,
often on logarithmic scales. \src{15.PNG}
\end{keydef}

\begin{formula}{Ranking by class}
\begin{tabular}{@{}lll@{}}
\toprule
\textbf{Property} & \textbf{Highest} & \textbf{Lowest} \\
\midrule
Density           & metals                & polymers \\
Elastic modulus   & metals, ceramics      & polymers \\
Tensile strength  & composites, metals, ceramics & none \\
\bottomrule
\end{tabular}

Composites and woods occupy a wide middle range for density and stiffness. For
tensile strength every class spans a range, so no class is lowest outright.
\src{15.PNG, 16.PNG, 17.PNG}
\end{formula}

\begin{keydef}{Ashby map (two-property optimization)}
Plots stiffness (elastic/Young's modulus, GPa) against density
($\mathrm{kg/m^3}$) on logarithmic axes, grouping materials into regions for
technical ceramics, non-technical ceramics, composites, metals, natural
materials, polymers, foams and elastomers. Used to compare and select materials
on two properties simultaneously. \src{18.PNG}
\end{keydef}

\topic{Materials Selection Process}

\begin{method}{Selecting a material and process}
\begin{steps}
\item For a specific application, determine the required properties:
mechanical, electrical, thermal, magnetic, optical, deteriorative.
\item From the list of required properties, identify candidate material(s).
\item For the best candidate material(s), specify processing technique(s), to
provide the required set of properties and to produce a component of the desired
shape and size. Example techniques: casting, mechanical forming, welding, heat
treating.
\end{steps}
\src{19.PNG}
\end{method}
% <<< APORE BODY END
\end{document}
